\documentclass{article}

\usepackage{full page}  % make the margins somewhat smaller than the default
\usepackage{pdfpages}  % needed to include pdf
\usepackage{hyperref}
\usepackage{color}
\usepackage{graphicx}

\usepackage{minted}



\title{Octave Keyboard}
\author{Daniel Chen and Vivian Hu}
\date{August 26, 2014}

\begin{document}

\maketitle
% \begin{figure}[ht!]
%   \centering
%   \includegraphics[width=2in, height=1.3in]{img/keyboard.jpg}
% \end{figure}

% abstract
% Give a concise summary of the goals and results from your project.
This report

\tableofcontents


\newpage
\section{Introduction and Problem Statement} % State the problem to be solved.
This report goes

This keyboard is a simple one­ octave keyboard with the capability to autoplay ``Kids'' by MGMT when a switch is turned on. When the keys are pressed, the appropriate notes are played through the speaker, and LEDs which correspond to the keys that are lit up. The LEDs can be disabled using a switch.





\section{Design Solution}

  \subsection{Specifications}
  
    % Specifications: Describe what your circuit does. This section will be similar to the project proposal. List any inputs such as buttons, keypads, sensors, etc., and any outputs such as displays, LEDs, actuators, etc. Include a block diagram of your circuit as well as a diagram of the control panel, showing the locations of buttons, displays, lights, etc., relative to each other. 


    \begin{description}

      \item[The Constructor] \hfill \\

      \item[getAllowableCombinations()] \hfill \\

      \item[main()] \hfill \\

    \end{description}

  \subsection{Operating Instructions}

    % Operating instructions: Provide any information needed to set up and run your circuit. This is the User’s Manual part of your report.

  \subsection{Theory of Operation}
  
    % Theory of operation: Discuss in detail how your circuit works. Start with a high-level overview of your design, illustrated by the block diagram. Give a prose description of each of the blocks and the interfaces between them. Describe what each block does and how it works. Make sure you discuss any tricks or functions that are not obvious from the logic diagram (e.g., ROMs used as lookup tables). From the information given in this section (with references to additional diagrams in the appendices), we should be able to reconstruct your design.

  \subsection{Construction and Debugging}

    % Construction and debugging: Describe how you built and debugged your circuit. Which components did you build first? How did you test them? What problems did you have?

    \begin{minted}{vhdl}
process
begin
  CLK <= '1'; wait for 10 NS;
  CLK <= '0'; wait for 10 NS;
end process;
    \end{minted}

\section{Evaluation of Design}
  
  % Justification and evaluation of your design: Why is your solution better or worse than some reasonable alternatives? What might you do differently if you had it to do over again?

\section{Conclusions and Recommendations}

  % Conclusions: Summarize the goals and accomplishments of your project. Give a comparison of your original proposal (and final specifications) and the degree to which you achieved your proposed objectives. What recommendations do you have for future groups considering such a project? What words of general advice would you give to future groups?

\section{Acknowledgments}

  % Acknowledgments: Technical reports and papers often have an acknowledgments section that thanks people other than the authors for their contributions. Use this section (if you wish) to acknowledge various sources (e.g. course staff, students, faculty, or others) for design ideas and other assistance.
  
  % In addition, you should use this section to explain the contributions of each partner to the project. (This is required.) For example, if each project partner was responsible for designing a different block of the system, state that here. If one partner had more responsibility for design, another for debugging, etc., describe that in this section. If different partners were responsible for different parts of the report, note that here.

\section{References}

  % References: Cite any references (books, application notes) used in your design (e.g., for some clever circuit design). It is not necessary to acknowledge parts specifications (data sheets), unless you are quoting.

\section{Appendices}

  % Appendices: All the diagrams, etc, you collected go here.

  \subsection{System level diagrams}

    \subsubsection{Front Panel}
      
      % Front panel: An annotated digital photo of your project (Nexys2 board + add-ons) is helpful to show how you use the buttons, switches, LEDs, etc, and also where PMods plug in, etc.  

    \subsubsection{Block Diagram}
      % A functional block diagram which gives an overview of the hardware functional modules. By looking at this diagram, anyone else who has had the course should be able to understand what your project does. It is not necessary to use a drawing program to create your block diagram. A piece of graph paper and a ruler work just fine.

      % Each block, each input / output port in a block, and each net connecting blocks should be labeled with a descriptive name. Input ports for a drawing are always at the left edge of the drawing, and output ports are at the right edge. Additionally, ports that connect to/from Pmods or the FX2 breadboard must be labeled with a connector and pin number, e.g., JA-2 for pin 2 of the JA Pmod connector.
      % Here is a suggested hierarchy for your block diagrams:
      %   (1) The top level that shows the main functional blocks, including switches, displays, etc. At this level, the logic in your FPGA can be one big “black box” with ports labeled to match the port declarations in the top level of your logic design (whether VHDL or schematic).
      %   (2) The top level of logic in your FPGA. It’s simple to base your drawing on the RTL schematic generated by the Xilinx tool. If necessary, annotate it by hand so the reader can easily see what each block does.
      %   (3) Work down “recursively” into each block. Make sure your drawings are labeled in such a way that their place in the hierarchy is clear.
    
    \subsubsection{Schematic Diagram}
      % A schematic diagram with signal names and pin numbers for any off-board circuitry that you designed and built (i.e., on the FX2 board). The purpose of this diagram is to show how your project is wired up and to help in debugging if it doesn’t work. Make sure you include detailed schematic diagrams for any special circuits (e.g. LED arrays). Each package should be given a unique identifier, known as the reference designator: Un for integrated circuits (where n = 1, 2, ...), Rn for resistors, Cn for capacitors, Qn for transistors, Sn for switches, Pn and Jn for plugs and jacks, respectively, Dn for diodes. LEDn is sometimes used for LEDs or arrays of LEDs (like 7- segment displays). If you have other parts, ask me.% 

    \subsubsection{Package Map}
      % A package map showing where each component is located on your board. This only applies to off-board circuits that you built. The package map should resemble an "aerial photograph" showing clearly all connectors (labeled Jn or Pn), sockets, ICs (Un), switches (Sn), displays, etc. The package map is used as an intermediate stage between the logic diagram and your circuit.

      % Suppose you are troubleshooting your project. By studying the symptoms and your logic diagram, you deduce that there may be a wiring error at a particular part, e.g., a keypad encoder. On the logic diagram, that chip is labeled U8. You find U8 on the package map and it shows you precisely where to look on your circuit board to find the suspicious part. An easy way to create a package map is to annotate a digital photo of the board.

    \subsubsection{Parts list}
      % Parts list: List all external components, including Pmods and parts that you wired up on the FX2 board. You don’t need a parts list if everything is confined to the Nexys board. Use sensible groupings (IC’s, discrete components, switches, connectors, etc.), and keep this list compact but complete. Omit trivial details like wire and wirewrap sockets.

  \subsection{Programmed Logic}

    \subsubsection{State Diagrams}
      % for all state machines

    \subsubsection{VHDL Code}
      % for each module

    \subsubsection{Resource utilization}
      % for FPGA

  \subsection{Memory Map}
    % Memory map: If you are using a memory chip or block memory in the FPGA as a complex look-up table or for data storage, show how the address and data bit patterns are partitioned among the logical variables involved (e.g., for storing sound clips, address bits 12-13 are which sound, and address bits 0- 11 are samples within a sound).

  \subsection{Timing Diagram}
    % A timing diagram, e.g., generated from a behavioral simulation, showing all important signals. It is not necessary to duplicate timing diagrams which may be found in a data sheet or user’s manual (e.g., for memory chips).


  % ADDITIONAL SUBSECTIONS?
  % 6. Copies of data sheets for any special components you used that aren’t already on the class website.
  % 7. Computer programs (e.g., Matlab codes used to generate lookup tables, etc).
  % 8. Other documentation. Include anything else that might be needed to understand or reproduce your design.


\end{document}


